Our second application of the above theorems is to a decades old problem posed by Stanley, who introduced and studied local $h$-vectors in the special case where $E = \emptyset$ and asked for a characterization of triangulations for which they vanish \cite[Problem~4.13]{Stanley92}. This problem remains open, and is of enduring interest \cite[Problem~2.12]{Athanasiadis16}. The extension to the case where $E$ is not empty is particularly relevant for applications to the monodromy conjecture \cite{Igusa78, DenefLoeser98, Stapledon17}. In \cite{LarsonPayneStapledon}, we prove a theorem on the structure of geometric triangulations with vanishing local $h$-vectors that is tailored to this purpose, and we use it to prove the monodromy conjectures for all singularities that are nondegenerate with respect to a simplicial Newton polyhedron.  See Theorems~1.1.1, 1.4.5, and 4.1.3 in loc. cit..

Here, we apply  Theorem~\ref{thm:resolution} to prove another theorem on the structure of faces in triangulations with vanishing local $h$-vectors.  Let $F \in \lk_\Gamma(E)$ be a face such that $F\sqcup E$ is interior. Following terminology from the monodromy conjecture literature (see, e.g. \cite{LemahieuVanProeyen11}), we say that $F$ is a \emph{pyramid with apex $w \in F$}  if $(F \sqcup E )  \smallsetminus w$ is not interior. Let $$\mathcal{A}_F  \coloneqq  \{ w \in F : F \mbox{ is a pyramid with apex } w \}, \mbox{ \ and \ } V_w  \coloneqq  \sigma( (F \sqcup E) \smallsetminus w)^c.$$
The elements of $V_w$ correspond to the \emph{base directions} of $F$, i.e. the facets of $2^V$ that contain the base of $F$, when viewed as a pyramid with apex $w$.  We say $F$ is a $U$-pyramid if there is an apex $w \in \mathcal{A}_F$ such that $|V_w|  = 1$. In other words, a $U$-pyramid is a pyramid with a unique base direction, for some choice of apex.

\begin{definition}
Let $F \in \lk_{\Gamma}(E)$ be a face.  An \emph{interior partition} of $F$ is a decomposition
\[
F = F_1 \sqcup F_2 \sqcup \mathcal{A}_F
\]
such that $F_1 \sqcup \mathcal{A}_F \sqcup E$ and $F_2 \sqcup \mathcal{A}_F \sqcup E$ are both interior.
\end{definition}

\begin{theorem}\label{thm:interiornonvanish}
Suppose $\ell(\Gamma,E) = 0$ and $F \in \lk_\Gamma(E)$ has an interior partition $F = F_1 \sqcup F_2 \sqcup \mathcal{A}_F$ such that $|F_1| \leq 2$.  Then $F$ is a $U$-pyramid.
\end{theorem}

See Remark~\ref{r:specialcase} for a short proof in a special case that illustrates the naturality of the 
$U$-pyramid condition.
The method of proof breaks down when $|F_i| \geq 3$.  See Example~\ref{example:nonrestriction}.  


\begin{remark}
The analogous theorem in \cite{LarsonPayneStapledon} requires that the triangulation be geometric and that the interior partition satisfies the additional condition $\sigma(F_2 \sqcup E)^c = \bigcup_{w \in \mathcal{A}_F} V_w$.  But then the hypothesis that $|F_1| \leq 2$ is dropped entirely.  So, even for geometric triangulations, there are cases of Theorem~\ref{thm:interiornonvanish} that are not necessarily covered by \cite[Theorem~4.1.3]{LarsonPayneStapledon}. It should be interesting to look for a common generalization of these vanishing results, and to pursue further progress on Stanley's problem of characterizing triangulations with vanishing local $h$-vector more generally.
\end{remark}


\begin{remark}
To the best of our knowledge, all of the theorems stated in the introduction are new even for regular triangulations. The reader who prefers to do so may safely restrict attention to geometric or even regular triangulations.  However, while the structure results for triangulations with vanishing local $h$-vectors in \cite{dMGPSS20} and \cite{LarsonPayneStapledon} rely on special properties of geometric triangulations, the proofs presented here work equally well for quasi-geometric homology triangulations, and we find it  natural to work in this level of generality.
\end{remark}
